% SPDX-License-Identifier: 0BSD
\section{Compiled, Eager, and Recurrence Execution}
\label{sec:eager-execution}

The three execution modes differ in how they turn the same current recursion
into numerical work. Compiled mode specializes groups of current relations to
one process. Eager mode applies reusable model kernels according to a
process-specific schedule. Recurrence mode constructs a compact schedule of
currents, contributions, and closures directly. These choices change the
balance between process-generation cost and runtime work, but not the
observable being evaluated.

Recurrence JIT O2 is the default execution mode.

\begin{center}
\small
\begin{tabular}{@{}L{0.72in}L{1.85in}L{2.00in}L{1.40in}@{}}
\toprule
\textbf{mode} & \textbf{process preparation} & \textbf{evaluation} &
  \textbf{characteristic balance} \\
\midrule
Compiled
  & Builds process-specific stage evaluators
  & Evaluates a small number of process-wide, multi-output stages
  & More preparation; greater fusion across a process \\
Eager
  & Maps the process graph to reusable prepared-model kernels
  & Applies those kernels according to the process dependencies
  & Reuses model work across processes \\
Recurrence
  & Builds a compact schedule directly from the current recursion
  & Evaluates scheduled currents, contributions, and closures
  & Avoids construction of a generic process DAG \\
\bottomrule
\end{tabular}
\end{center}

\subsection{Common physics contract}

All three modes retain the same external states, perturbative orders,
helicities, colour components, coherent amplitude groups, symmetry factors,
and initial-state averages. In a contracted colour basis their common
observable has the schematic form
\begin{equation}
  \mathcal M_p
  =\mathcal N\sum_h\sum_{a,b}
    A_{ha}(p)^{*}\,C_{ab}\,A_{hb}(p),
  \label{eq:three-mode-common-observable}
\end{equation}
with the corresponding resolved form retained when helicity or leading-colour
flow components are requested. Each mode must reproduce both the directly
evaluated total and the sum of those resolved components.

Current contributions are accumulated coherently before a propagator or other
current transformation is applied. Final closures then form the amplitude
components entering Eq.~\eqref{eq:three-mode-common-observable}. This ordering
is common to the three modes and prevents execution strategy from changing the
normalization.

\subsection{Prepared-model boundary}

A prepared model contains reusable numerical information derived from the
model's particles, parameters, propagators, and interaction tensors. Eager and
recurrence JIT O2 measurements reuse this model-level preparation and add only
the process-dependent schedule. Compiled JIT O3 measurements instead include
construction and compilation of their process-specific evaluators.

Prepared-model creation occurs outside the process-generation timer. The
reported eager and recurrence generation times therefore describe adding a
process to an already prepared model, not importing and compiling the model
itself.

\subsection{Batches, selectors, and portability}

Evaluation is batched over phase-space points, allowing the JIT backend to use
the point dimension for vectorization. Partitioning a large input batch changes
neither the process representation nor the numerical result.

Helicity and colour-flow selectors refer to physical resolved components. A
complete leading-colour calculation retains every declared helicity and flow,
so selecting one flow with summed helicities or one helicity with all flows
does not alter the underlying physics coverage. The two layouts used for these
workloads are described in Sec.~\ref{sec:lc-flow-layouts}.

Portable JIT state is compiled for the receiving processor when loaded.
Target-native ASM and C++ payloads remain architecture-specific, and the exact
symbolic representation is retained where higher-precision evaluation is
supported.

\subsection{Interpretation}

Compiled mode may benefit most when process-wide fusion outweighs its
additional preparation cost. Eager mode can reduce repeated preparation when
many processes share a model, while recurrence provides a compact direct
representation of the current construction. The tables compare these
trade-offs under identical physics, selector, precision, and validation
conditions; no mode is assumed to be universally fastest.
